\documentclass[11pt]{article}

\usepackage[margin=1in]{geometry}
\usepackage{booktabs}
\usepackage{array}
\usepackage{longtable}
\usepackage{xcolor}
\usepackage{hyperref}
\usepackage{amsmath}
\usepackage{enumitem}
\usepackage{titlesec}
\usepackage{fancyhdr}
\usepackage{caption}

\definecolor{verdictred}{RGB}{150,30,30}
\definecolor{verdictgreen}{RGB}{30,110,30}
\definecolor{headergray}{RGB}{90,90,90}

\hypersetup{
  colorlinks=true,
  linkcolor=black,
  citecolor=black,
  urlcolor={blue!60!black},
  pdftitle={PEARL-Neuro EEG Classification -- Final Report},
  pdfauthor={PEARL-Neuro Analysis Team}
}

\titleformat{\section}{\large\bfseries}{\thesection}{0.6em}{}
\titleformat{\subsection}{\normalsize\bfseries}{\thesubsection}{0.6em}{}

\pagestyle{fancy}
\fancyhf{}
\fancyhead[L]{\small\color{headergray} PEARL-Neuro EEG Classification -- Final Report}
\fancyhead[R]{\small\color{headergray} \thepage}
\renewcommand{\headrulewidth}{0.4pt}

\newcommand{\verdict}[2]{%
  \noindent\fbox{\parbox{\dimexpr\linewidth-2\fboxsep-2\fboxrule}{%
  \textbf{\color{#1}#2}}}%
}

\title{\vspace{-1.5cm}PEARL-Neuro EEG Classification\\[4pt]
\large Predicting Genetic Alzheimer's-Risk Group from Resting-State and Task EEG\\[2pt]
\large Final Report}
\author{Dataset: OpenNeuro \texttt{ds004796} \\ 79 EEG-imaged, genotyped subjects (of 192 genotyped total)}
\date{Prepared 2026-08-13}

\begin{document}
\maketitle
\thispagestyle{fancy}

\begin{abstract}
\noindent
This report documents two attempts to classify genetic Alzheimer's-risk
group (APOE~$\varepsilon4$ + PICALM status) from EEG, using a
pitch-synchronous waveform-shape feature family (PSR/PSWT) evaluated
against a conventional spectral baseline. The original analysis
(Phases~1--4) produced a validated \textbf{null result} (AUC 0.474,
$p=0.577$). A second attempt (Phase~5) fixed two real preprocessing bugs
and obtained a marginal \textbf{positive result} (AUC 0.651, $p=0.041$).
That positive result was then subjected to three independent validation
checks -- resampling stability, single-fix attribution, and full-precision
reproducibility -- and \textbf{did not survive}: it is statistically
indistinguishable from a covariates-only baseline once resampled, and
neither underlying preprocessing fix reproduces it in isolation. A final
phase then \textbf{bounded} the null: on this same cohort and pipeline, a
representation that successfully decodes a control target (sex, AUC 0.688,
$p=0.019$) finds nothing for genotype (AUC 0.506, $p=0.48$), establishing
that the measurement is sensitive enough to detect a real subject-level
effect and that no genotype effect is present at this sample size. The
honest conclusion of this engagement is that no reliable EEG-based
classifier for this target emerged from this dataset, that the null is
now bounded rather than merely inconclusive, and that the dataset's own
hard cap (79 subjects, a detection floor of AUC 0.693 at 80\% power, no
independent replication cohort) makes this the practical ceiling of what
can currently be known -- not a shortfall of the analysis.
\end{abstract}

\vspace{0.3cm}
\verdict{verdictred}{Final verdict: NOT VALIDATED. Do not deliver AUC 0.651 as a positive finding. The null is bounded, not inconclusive -- see \S 5 and \S 6.}
\vspace{0.4cm}

\tableofcontents

% =====================================================================
\section{Objective and Background}

The client requested a genetic Alzheimer's-risk classifier built from EEG
recordings in the PEARL-Neuro dataset (\texttt{ds004796}), comparing a
custom pitch-synchronous waveform-shape feature family (PSR/PSWT) against
a conventional spectral baseline, benchmarked against a published result
of AUC~$\approx 0.58$ (Li et al., 2025, \textit{Sensors} 25(1):52, MSIT
task).

The dataset provides EEG and genotype data for 192 participants; only 79
have both EEG recordings and confirmed genotype status, and this project
uses all 79. Genetic risk group (\texttt{N} / \texttt{A\_P\_minus} /
\texttt{A\_P\_plus}) is derived from APOE~$\varepsilon4$ carrier status
combined with PICALM genotype, per the client's group-assignment
convention (independently reconciled against official genotype columns in
Phase~0).

The engagement proceeded in two parts:

\begin{enumerate}[label=\textbf{Part \arabic*.}, leftmargin=2.2cm]
  \item \textbf{Phases 0--4} (original engagement): audit, preprocessing,
        feature engineering, and modeling under a fully frozen,
        label-blind protocol, ending in a validated null result.
  \item \textbf{Phase 5} (client-authorized second attempt): a genuine
        attempt to close the gap to the published benchmark by fixing
        two preprocessing issues surfaced during Phase~1--4, followed by
        rigorous validation of whatever result emerged.
\end{enumerate}

% =====================================================================
\section{Data and Methods}

\subsection{Dataset and cohort}

\begin{itemize}[leftmargin=1.2cm]
  \item \textbf{Source:} OpenNeuro \texttt{ds004796}, resting-state EEG
        plus two task paradigms (MSIT, Sternberg working-memory task).
  \item \textbf{Population:} 79 subjects with both EEG and confirmed
        genotype. After preprocessing QC and the frozen exclusion rules,
        the analysis cohort used in the primary comparison is
        $n=64$ subjects in every reported condition.
  \item \textbf{Primary target:} \texttt{binary\_risk\_vs\_none}
        (at-risk genotype group vs.\ no-risk group) -- declared before any
        feature existed, and never changed despite cohort size moving
        across phases.
\end{itemize}

\subsection{Preprocessing}

Standard EEG preprocessing pipeline: montage assignment, average
referencing, band-pass filtering (0.5\,Hz high-pass, explicit anti-alias
low-pass), 50\,Hz line-noise removal (Polish mains), automated bad-channel
detection and interpolation, resampling, and ICA-based artifact removal
(ICLabel-classified components: muscle, eye blink, heartbeat, line noise,
channel noise).

Every preprocessing threshold is fixed in a frozen configuration file and
applied uniformly \emph{before} any genotype label is read (label-blind
QC), enforced by dedicated tests that scan the codebase for any import
path to label-bearing data outside two narrowly-scoped, audited reporting
modules.

\paragraph{Two bugs found and fixed in Phase 5.} Reviewing the original
pipeline surfaced two real, independently-justified issues:

\begin{enumerate}[leftmargin=1.2cm]
  \item \textbf{Reference-order bug.} The average EEG reference was
        computed \emph{before} bad-channel detection and interpolation,
        letting a noisy channel contaminate the reference signal used by
        every other channel. Fixed by reordering the pipeline so
        interpolation happens first.
  \item \textbf{ICA component-count re-evaluation.} The original pipeline
        used a fixed number of ICA components (30) for every subject.
        Comparison data showed this caused excessive component removal on
        subjects with more bad channels. Replaced with explained-variance
        based component selection (0.99), scaled to each subject's actual
        post-interpolation channel rank.
\end{enumerate}

\subsection{Feature engineering}

Two feature families were extracted from the resting-state eyes-closed
window, per subject:

\begin{itemize}[leftmargin=1.2cm]
  \item \textbf{PSWT (primary):} pitch-synchronous waveform-shape
        features per detected alpha-band cycle -- harmonic amplitude
        profile (2nd--5th harmonics), harmonic-to-interharmonic ratio,
        cycle-to-cycle variability, period variability, rise-decay ratio,
        peak-trough sharpness ratio (18 features total).
  \item \textbf{Baseline spectral (comparison):} conventional
        band-power/alpha-peak features over the same channels and window
        (23 features total), used as the point of comparison against the
        published benchmark's spectral approach.
\end{itemize}

A richer multitaper time-frequency-area (TFA) feature family was also
built and tested in Phase~5 \emph{only} against the benchmark-reproduction
question (never against the primary target directly, to avoid conflating
two separate questions in one modeling attempt); it did not close the
benchmark gap (AUC 0.546 vs.\ 0.459 prior, $p=0.315$, not significant).

\subsection{Modeling and evaluation protocol}

A single protocol, declared before any label was joined and never changed
after seeing a result:

\begin{itemize}[leftmargin=1.2cm]
  \item \textbf{Model:} L2-regularized logistic regression (declared as
        primary before any result existed; tree ensembles, gradient
        boosting, and neural networks were explicitly excluded from the
        primary analysis).
  \item \textbf{Cross-validation:} 5-fold $\times$ 10 repeats,
        group-stratified, pooled out-of-fold AUC across repeats.
  \item \textbf{Inference:} full-pipeline permutation test (1000
        permutations) for the $p$-value, subject-level bootstrap (2000
        resamples) for the 95\% confidence interval.
  \item \textbf{Reference lines:} chance (0.5), a QC-metrics-only
        classifier (contamination-floor diagnostic, not a result), a
        nuisance-only classifier (age/sex/education/BDI/SES + QC
        covariates, no EEG features), and the published benchmark (0.58).
  \item \textbf{Stopping rule:} a maximum of three feature families,
        declared up front, to prevent post-hoc family-shopping.
\end{itemize}

% =====================================================================
\section{Results: Original Analysis (Phases 1--4)}

\begin{table}[h]
\centering
\caption{Phase 3 primary result.}
\begin{tabular}{lc}
\toprule
Metric & Value \\
\midrule
AUC & 0.474 \\
95\% CI & [0.266, 0.702] \\
Permutation $p$ & 0.577 \\
Cohort & $n=64$ (39 no-risk / 25 at-risk) \\
Verdict & \textbf{NULL} \\
\bottomrule
\end{tabular}
\end{table}

Two positive controls were run, both on the Stage~5 corrected-preprocessing
cohort (the original, bug-present preprocessing scored 0.830/0.486/0.459 on
the same three checks, respectively -- consistently close, so preprocessing
is not the source of the pattern below). \textbf{Eyes-open vs.\
eyes-closed} (within-subject, strong control) passed clearly: AUC 0.849
against a 0.80 floor -- the features correctly measure alpha waveform
shape and generalise \emph{within} a subject. \textbf{Sex classification}
(cross-subject control) \emph{failed}: AUC 0.465 against a declared 0.65
floor -- sex is reliably decodable from resting EEG in adequately powered
samples, so this failure means the feature representation, as built, does
not demonstrably generalise \emph{across} subjects for a genuinely
decodable trait. This does not mean the pipeline is broken (1a passing
rules that out), but it does mean the genotype null cannot be
distinguished from ``this representation discards the between-subject
information needed to detect any cross-subject effect'' -- a real
limitation of what has been validated, not just of what was found.
Benchmark reproduction on the MSIT task landed at chance (AUC 0.490), a
characterized but unresolved gap attributed to reference-order and ICA
issues (addressed in Phase~5) and a feature-richness gap.

% =====================================================================
\section{Results: Phase 5 Attempt and Its Validation}

\subsection{The positive result}

Re-running Phase 3's exact frozen methodology on Phase 5's corrected
preprocessing (both fixes applied) produced a verdict flip:

\begin{table}[h]
\centering
\caption{Original vs.\ Phase 5 (both preprocessing fixes applied).}
\begin{tabular}{lcc}
\toprule
 & Original (Phase 3) & Phase 5 (corrected preprocessing) \\
\midrule
AUC & 0.474 & 0.651 \\
95\% CI & [0.266, 0.702] & [0.353, 0.764] \\
Permutation $p$ & 0.577 & 0.041 \\
Verdict & NULL & POSITIVE (marginal) \\
Cohort & $n=64$ (39/25) & $n=64$ (38/26), 5 subjects different \\
\bottomrule
\end{tabular}
\end{table}

Four reasons for caution were flagged immediately, before any further
validation work: $p=0.041$ is barely under the 0.05 threshold; the 95\%
CI's lower bound (0.353) approaches chance; the point estimate moved by
0.177~AUC from just 5 of 64 subjects changing (a side effect of corrected
QC, not a deliberate choice); and the result did not corroborate itself
across related secondary comparisons (none survived correction for
multiple comparisons).

\subsection{Validation of the positive result}

Three further, independent checks were run before any delivery decision.

\paragraph{Stability under resampling.} The same 64-subject sample was
resampled 325 times (175 leave-$k$-out draws, $k=2$--8; 150
subject-level bootstrap resamples) and the pooled CV AUC recomputed each
time.

\begin{table}[h]
\centering
\caption{Resampling stability of the 0.651 point estimate.}
\begin{tabular}{lccccc}
\toprule
Method & $n$ & 5th pct. & median & 95th pct. & \% below chance \\
\midrule
Leave-$k$-out & 175 & 0.531 & 0.625 & 0.715 & 1.1\% \\
Bootstrap & 150 & 0.407 & 0.565 & 0.726 & \textbf{24.0\%} \\
\bottomrule
\end{tabular}
\end{table}

The bootstrap median (0.565) sits almost exactly on the nuisance-only
reference line (0.557) -- the resampling distribution's center shows no
meaningful separation from ``covariates alone,'' and a full quarter of
resamples fall below chance entirely. This is not the signature of a
reliably-located effect.

\paragraph{Single-fix attribution.} Stage 1's preprocessing changed two
things at once. Each was isolated in a separate full preprocessing
re-run, holding the other fix at its default.

\begin{table}[h]
\centering
\caption{Full picture across every condition tested.}
\begin{tabular}{lcccc}
\toprule
Condition & AUC & 95\% CI & $p$ & Significant? \\
\midrule
Original (both bugs present) & 0.474 & [0.266, 0.702] & 0.577 & No \\
Reference-order fix alone & 0.583 & [0.324, 0.750] & 0.234 & No \\
ICA fix alone & 0.602 & [0.315, 0.739] & 0.174 & No \\
\textbf{Both fixes together} & \textbf{0.651} & \textbf{[0.353, 0.764]} & \textbf{0.041} & \textbf{Yes (marginal)} \\
\bottomrule
\end{tabular}
\end{table}

\textbf{Neither fix, applied alone, reproduces significance.} Only the
specific combination, on this specific 64-subject sample, crosses
$p<0.05$ -- a pattern more consistent with a sample-specific coincidence
than with two real, additive biological signals.

\paragraph{Reproducibility.} After restoring the production preprocessing
tree to the both-fixes-applied state, the primary analysis was re-run at
full statistical precision: AUC 0.6509, CI [0.353, 0.764], $p=0.0410$ --
matching the original 0.651 to within rounding. This confirms the number
was computed correctly; it does not, on its own, make the number
reliable -- that question is answered by the two checks above, both of
which argue against it.

% =====================================================================
\section{Phase 6: Bounding the Null}

A null result carries an inherent ambiguity: did the analysis find nothing
because there is nothing to find, or because the measurement was insensitive?
Phase 3's control battery left exactly this open. The pitch-synchronous feature
family scored \textbf{AUC 0.465 on sex} -- a \emph{failed} positive control
against a declared floor of 0.65. Since sex is reliably decodable from resting
EEG in adequately powered samples, that failure meant the genotype null could not
be distinguished from a representation that simply discards all between-subject
information.

Phase 6 resolved this by testing four already-computed representations against
the sex control, on the identical cohort, preprocessing, and cross-validation
protocol used for the genotype analysis. \textbf{No new genotype analysis was
performed.}

\begin{table}[h]
\centering
\caption{Sex control across four representations ($n=64$, full precision).}
\begin{tabular}{lccccc}
\toprule
Representation & Cols & AUC & 95\% CI & perm $p$ & Verdict \\
\midrule
\texttt{baseline\_rest} & 9 & \textbf{0.688} & [0.431, 0.816] & \textbf{0.019} & PASS \\
\texttt{baseline\_rest\_msit} & 16 & \textbf{0.696} & [0.421, 0.826] & \textbf{0.007} & PASS \\
\texttt{baseline\_all} & 23 & \textbf{0.652} & [0.373, 0.819] & \textbf{0.029} & PASS \\
\texttt{pswt} & 18 & 0.465 & [0.272, 0.700] & 0.619 & FAIL \\
\bottomrule
\end{tabular}
\end{table}

\textbf{Three of the four representations decode sex significantly.} Only the
pitch-synchronous family fails, and it fails decisively ($p=0.62$, essentially
chance). The cohort is not inert -- it carries subject-level information that a
suitable representation can recover.

\subsection{The bounded null}

The decisive comparison holds the representation fixed and varies only the
target. The same nine resting-state spectral features:

\begin{table}[h]
\centering
\begin{tabular}{lccc}
\toprule
Target & AUC & $p$ & Outcome \\
\midrule
Sex (control) & \textbf{0.688} & \textbf{0.019} & Signal detected \\
Genotype (primary) & 0.506 & 0.482 & Nothing \\
\bottomrule
\end{tabular}
\end{table}

Matched on features, cohort, preprocessing, and protocol, this representation
detects a real subject-level trait and detects nothing for genotype.
\textbf{The genotype null is therefore a statement about genotype, not an
artifact of an insensitive measurement.} This upgrades the finding from
\emph{inconclusive} to \emph{bounded by a demonstrated negative}.

Two pre-declared verification checks passed: the \texttt{pswt} control
reproduced the earlier result bit-exactly (0.4645 / [0.2715, 0.7000] /
$p=0.6194$), confirming the pipeline is deterministic and the stored derivatives
are intact; and \texttt{baseline\_rest\_msit} reproduced an earlier
16-column result to within 0.003 across a preprocessing change.

\subsection{The detection ceiling, quantified}

A formal power analysis for this cohort ($n=34$ vs 30, $\alpha=0.05$) establishes
what it can and cannot resolve:

\begin{table}[h]
\centering
\begin{tabular}{lc}
\toprule
Quantity & Value \\
\midrule
Minimum detectable AUC at 80\% power & \textbf{0.693} \\
Minimum detectable AUC at 90\% power & 0.720 \\
Power to detect the published benchmark (0.58) & \textbf{20\%} \\
Power to detect AUC 0.65 & 57\% \\
\bottomrule
\end{tabular}
\end{table}

All three passing sex results sit at or near that 80\%-power floor (0.688,
0.696, 0.652 against an MDE of 0.693). The cohort can \emph{just barely} resolve
effects of that magnitude. Sex is among the most robustly encoded subject-level
traits in resting EEG; genotype risk group is substantially subtler. This is the
quantitative basis for the recommendation against further analysis of this
cohort: at 20\% power against the benchmark effect, a study would fail to detect
it four times in five even if it were real and exactly as large as published.

\textbf{Limitation of this bound.} It applies to the representations tested. It
does not prove that no representation could recover genotype signal -- only that
the ones examined here, including one demonstrably capable of subject-level
decoding, do not.

% =====================================================================
\section{Honest Evaluation}

\verdict{verdictred}{The Stage 3 positive result (AUC 0.651) does not survive validation and should not be delivered as a positive finding.}

Two independent methods -- resampling stability and single-fix
attribution -- reach the same conclusion by different routes, which is
stronger evidence than either alone. This pattern matches a well-known
statistical phenomenon: underpowered studies that happen to cross a
significance threshold systematically overstate the true effect (the
``winner's curse''; Button et al., 2013, \textit{Nature Reviews
Neuroscience}). A formal power analysis independently confirms this
study's minimum detectable effect at 80\% power (Cohen's $d=0.596$) is
\emph{larger} than the effect actually observed ($d=0.549$) -- this
sample is below its own floor for reliable detection, and power to detect
an effect the size of the published 0.58 benchmark was only 27\%.

\subsection{What is, and is not, being claimed}

\begin{itemize}[leftmargin=1.2cm]
  \item \textbf{Not claimed:} that PEARL-Neuro EEG data contains no
        genotype-related signal whatsoever. Absence of a validated
        finding at this sample size is not proof of absence of an
        effect.
  \item \textbf{Claimed, with evidence:} that this 79-subject,
        single-cohort dataset cannot currently support a validated
        positive claim for this target, and that the one candidate
        positive result obtained does not survive the standard checks a
        credible finding must pass.
  \item \textbf{Newly established (Phase 6):} the null is
        \emph{bounded}, not merely inconclusive. A representation that
        demonstrably carries subject-level information on this exact
        cohort (sex, AUC 0.688, $p=0.019$) finds nothing for genotype
        (AUC 0.506, $p=0.48$). The measurement is sensitive enough to
        detect a real subject-level trait; it detects no genotype
        effect.
  \item \textbf{Two real engineering improvements were made} regardless
        of the modeling outcome: the reference-order bug and the ICA
        component-count issue are both independently justified fixes,
        documented and tested, and are now the permanent default
        configuration for this pipeline going forward.
\end{itemize}

% =====================================================================
\section{Limitations}

\begin{itemize}[leftmargin=1.2cm]
  \item \textbf{Hard sample-size ceiling.} \texttt{ds004796} has exactly
        79 EEG-imaged, genotyped subjects; all are already used. There is
        no larger cohort to draw from within this dataset.
  \item \textbf{No independent replication cohort exists.} This dataset
        is, to available knowledge, the only public EEG dataset genotyped
        for APOE/PICALM Alzheimer's-risk status at this scale -- true
        independent replication would require new primary data
        collection, outside this engagement's scope.
  \item \textbf{Rest-vs-task benchmark mismatch.} The published 0.58
        benchmark used task-state EEG (MSIT), not resting-state; this
        project's own MSIT reproduction landed at chance, a
        characterized but unresolved gap, not a validated comparison
        point in either direction.
  \item \textbf{Cohort sensitivity.} A change of only 5 subjects (a QC
        side effect, not a deliberate resampling choice) moved the point
        estimate by 0.177 AUC -- this small a sample is inherently
        sensitive to exactly which subjects are included.
\end{itemize}

% =====================================================================
\section{Recommendations}

\begin{enumerate}[leftmargin=1.2cm]
  \item \textbf{Do not deliver AUC 0.651 as a positive finding} to any
        downstream audience (publication, product decision, funding
        pitch) without the full validation context in this report.
  \item \textbf{Deliver the honest characterization} of what was tried,
        what changed, why it looked promising initially, and why it does
        not survive scrutiny -- this is a complete and useful answer to
        the question actually asked.
  \item \textbf{If a definitive answer is required}, the only path
        forward is independent data collection: a new, adequately
        powered EEG + genotype cohort, sized using the power analysis in
        this project (well beyond 79 subjects for an effect in the
        0.55--0.65 AUC range).
  \item \textbf{Retain the two preprocessing fixes} (reference-order,
        ICA component count) as permanent pipeline improvements
        independent of the modeling outcome -- they are correct on their
        own merits.
\end{enumerate}

% =====================================================================
\section{Conclusion}

Fixing two real preprocessing bugs changed this project's primary result
from a clean null to a marginal, initially promising positive. Subjecting
that positive result to the same rigor that made the original null
trustworthy -- resampling stability, single-variable attribution, and
independent reproducibility checks -- showed it does not hold up.

A final phase then closed the remaining ambiguity. A null is only
interpretable if the measurement was capable of detecting something in the
first place, and that had never been established. It now has been: on this
exact cohort and pipeline, a representation that successfully decodes a
control trait (sex, AUC 0.688, $p=0.019$) finds nothing whatsoever for
genotype (AUC 0.506). The instrument works; the signal is not there.

The defensible deliverable from this engagement is not a discovery, but a
methodologically rigorous, independently corroborated answer: at 79
subjects, in a single cohort, with no replication data available, this
dataset does not support a validated EEG-based classifier for genetic
Alzheimer's-risk group -- and its detection floor (AUC 0.693 at 80\%
power, 20\% power against the published benchmark) explains why. That is a
complete and honest answer to the question this engagement set out to
resolve, and a quantitative specification of what a future study would
need to do better.

\vspace{0.5cm}
\hrule
\vspace{0.3cm}
{\small\color{headergray}
\noindent\textbf{Full artifact trail:} \texttt{pearl/reports/phase5\_final\_report.md},
\texttt{pearl/reports/phase5\_stage3\_validation.md},
\texttt{pearl/reports/phase5\_stage3\_stability.md},
\texttt{pearl/reports/phase5\_stage1b\_reference\_order\_only.md},
\texttt{pearl/reports/phase5\_stage1c\_ica\_only.md}, and
\texttt{HANDOVER.md} in the project repository.
}

\end{document}
